\section{Introduction}\label{sec:intro}
An LLM Agent receives a support ticket, checks an order and submits a refund. An observer of the encrypted path to the commerce service sees contacted endpoints, exchange sizes and waiting times. Recent Agent studies use such metadata to infer prompts, tasks, workflows and user traits~\cite{jeong2026wra,dangol2026workflow,pouryousef2026traffic,li2026ant}. Adding decoy calls can obscure the requested action, but each extra exchange must be generated somewhere and terminated somewhere. These choices determine its traffic signature and business cost.

\paragraph{The mechanism}
In the planner loop we measure, an ordinary tool call follows a model decision. A controller can insert a decoy directly after a tool response, bypassing the next model wait. That shortcut separates the decoy by its inter-call gap. Delaying the call changes this feature, but leaves its destination, position, size and relationship to the other calls available to an attacker. A complete shadow task changes those relationships as well as adding a second planner's work. We compare these implemented choices on actual TLS records and track the work performed with each trace.

Classical query obfuscation already links decoy quality to its generator~\cite{howe2008trackmenot,peddinti2010trackmenot,balsa2012obpws}. Our question is how that principle interacts with an Agent's model decisions and tool authority. Immediate cover is therefore a calibration arm, not a proposed defense: cover that ignores the honest timing law is known to be separable~\cite{shmatikov2006adaptive,peddinti2010trackmenot}, and we include it to confirm that the attack suite detects this in a live Agent loop. A byte-only attack can miss the controller's timing shortcut; a timing-only explanation can miss remaining call structure; a goal-only check can miss collateral business effects. The measurements therefore pair observation information with inference, task outcome, business state and resource use.

\begin{figure*}[t]
\centering
\begin{tikzpicture}[x=1cm,y=1cm,font=\small,>=Latex]
\tikzset{box/.style={draw,rounded corners,align=center,minimum height=.85cm,text width=2.4cm},
 note/.style={align=center,font=\footnotesize,text width=3.1cm}}
\node[box,fill=blue!5] (u) at (1.3,1.15) {User\\``Process the next ticket''};
\node[box,fill=blue!5] (a) at (5.1,1.15) {Agent\\plan, check, act, retry};
\node[box,fill=green!6] (m) at (9.2,2.6) {Model service\\next action / streamed output};
\node[box,fill=green!6] (t) at (9.2,-.35) {Tool service\\ticket, order, refund};
\node[box,fill=orange!8] (d) at (14.1,1.15) {Specialist Agent\\optional delegation};
\draw[<->,thick] (u)--node[above,yshift=15pt]{client link C}(a);
\draw[<->,thick] (a.north east)--node[above,sloped]{model link M}(m.west);
\draw[<->,thick] (a.south east)--node[below,sloped]{tool link T}(t.west);
\draw[<->,thick] (a.east)--node[above]{coordination A}(d.west);
\node[note] at (2.0,-.35) {\textbf{Transit observer}\\lengths, gaps, connections;\\destinations when available};
\node[note] at (13.7,-.35) {\textbf{Recipient / active service}\\received arguments;\\optionally selected responses};
\node[draw=black!45,fill=black!3,align=left,text width=15.7cm,inner sep=6pt] at (7.9,-1.65) {
\textbf{Joint protection decision:} Which observed traces remain distinguishable? Does cover preserve the requested goal, allowed business state and resource budget? State where the cover stops and whether its feedback reaches the planner.};
\end{tikzpicture}
\caption{Running support-ticket example and observation rights. A tool result returns to the Agent and can enter the next model request. Each encrypted link exposes its own traffic shape; a recipient additionally sees the request addressed to it. Concurrent tickets may share a connection on either link. The arrows represent application communication, not a one-to-one mapping from tool events to packets.}
\label{fig:chain}
\end{figure*}


\paragraph{Setting}
We compare four cover policies and no cover in 3,000 live executions of Qwen3:8B and Llama3.1:8B over MCP and TLS~1.3/HTTP/2 on one host. The provider cooperates by padding responses and permitting isolated execution; history-based policies also use server-side minimum holds. Section~\ref{sec:design} states these capabilities. When only request shape is private, protocol padding may suffice. Hiding response activity calls for additional responses, which a cooperating provider can synthesize or an application call can produce. Reaching the application need not require a business write.

\paragraph{Findings}
\begin{itemize}
\item \textbf{Time supports recognition even after sampled waits.} Immediate insertion exposes a model-wait shortcut. After sampling waits, an exploratory feature ablation still finds lower recognition without time in three of four history-dispatch cells after Holm correction. The concurrent comparisons fall from 99 to 66/120 (Qwen) and 106 to 77/120 (Llama). Qwen sequential remains nonsignificant after this family's correction ($p=.0931$). Explicit-position ablations show no significant difference in their separate family.
\item \textbf{Context decoding recovers recognition from associated encrypted records.} With test call windows withheld, context removal and joint decoding improve over refitted baselines in five of eight exploratory comparisons after Holm correction. Concurrent recognition rises from 56 to 80/120 and 78 to 112/120. The original short-gap remover underperformed on this view; the stronger results show why an unsuccessful removal pipeline cannot stand in for a privacy assessment.
\item \textbf{Shadow cover changes a complete policy's inference--cost trade-off.} The strongest tested call-window attack scores 70/120 and 64/120 at about twice the model tokens; record-only maxima are 73 and 69/120. Added context decoders do not improve the selected baseline in either shadow cell. The second business task changes call count, composition and terminal coverage together, so the comparison does not isolate online planning. Task, state and resource outcomes accompany those same executions; a separate stress control illustrates collateral effects missed by goal completion.
\end{itemize}

\paragraph{Contributions}
We contribute (1) live measurements and feature ablations that distinguish the model-wait shortcut from residual temporal context; (2) context and joint-decoding attacks on records without test call windows, alongside complete-policy inference, task, state and resource outcomes; and (3) an execution-aware evaluation method that states cover termination, state and feedback authority, and service support. The new analyses reuse the measured executions and are exploratory. Section~\ref{sec:floor} provides an ideal candidate-label reference for interpreting shadows, and Appendix~\ref{app:map} supplies the forty-four-setting background.
